\documentclass[11pt]{article}
\usepackage[T1]{fontenc}
\usepackage{amsmath,amssymb,amsthm}
\usepackage[margin=1in]{geometry}
\usepackage{hyperref}
\newtheorem{theorem}{Theorem}
\newtheorem{lemma}[theorem]{Lemma}
\newtheorem{conjecture}[theorem]{Conjecture}
\newtheorem{fact}[theorem]{Fact}
\title{Disproof of TLMC Conjecture 00000000119\\
\large The Pisano period $\pi(p)$ of a prime $p$ is prime only for $p=2$}
\author{TLMC submission 00000000119}
\date{September 2026}
\begin{document}
\maketitle

\begin{conjecture}[TLMC 00000000119]
There are infinitely many primes $p$ for which the Pisano period $\pi(p)$
(the period of the Fibonacci sequence $F_0=0,\,F_1=1,\,F_{n+2}=F_{n+1}+F_n$
modulo $p$) is itself prime.
\end{conjecture}

\section{Verdict and disproof}

The conjecture is \textbf{false}.  The set of primes whose Pisano period is
prime is exactly $\{2\}$:  the referee's own data gives $\pi(2)=3$ (prime)
and $\pi(3)=8$, $\pi(5)=20$, $\pi(7)=16$ (all even), and the obstruction
generalizes to every prime $p\ge 3$.

\begin{fact}[Matrix setup]
Let $Q=\begin{pmatrix}1&1\\1&0\end{pmatrix}$, so $\det Q=-1$ and
$Q^m=\begin{pmatrix}F_{m+1}&F_m\\F_m&F_{m-1}\end{pmatrix}$.
\end{fact}

\begin{lemma}[Parity obstruction]\label{lem:parity}
Let $p\ge 3$ and let $n\ge 1$ satisfy $F_n\equiv 0\pmod p$ and
$F_{n+1}\equiv 1\pmod p$.  Then $n$ is even and $n\ne 2$.
\end{lemma}

\begin{proof}
The congruences say $Q^n\equiv I\pmod p$, hence
$\det(Q^n)=(-1)^n\equiv 1\pmod p$.  Since $p\ge 3$, this forces $n$ even.
Moreover $n=2$ is impossible: $F_2=1\not\equiv 0\pmod p$.
\end{proof}

\begin{lemma}[Cassini]\label{lem:cassini}
For all $n$, $F_{n+1}^2-F_nF_{n+2}=(-1)^n$.
\end{lemma}

Equivalently, modulo an odd $p$ the period congruences of
Lemma~\ref{lem:parity} give $1\equiv(-1)^n$, which is absurd for odd $n$;
this identity is the engine behind the parity statement and is proved in
the Lean development (`\texttt{cassini\_even}`, `\texttt{cassini\_odd}`).

\begin{theorem}\label{thm:main}
The Pisano period $\pi(p)$ of a prime $p$ is prime if and only if $p=2$,
in which case $\pi(2)=3$.
\end{theorem}

\begin{proof}
$\pi(2)=3$ is prime (direct computation, `pisano2' in the Lean file).
Conversely let $p\ge 3$ be prime.  By Lemma~\ref{lem:parity} the period
$\pi(p)$ is even and different from $2$, hence composite.  So no prime
$p\ge 3$ has prime Pisano period, and the set of the conjecture is exactly
$\{2\}$.
\end{proof}

\section{Formal verification}

The Lean development (directory \texttt{lean4/}, core Lean~4 only, no
Mathlib) contains a fully axiom-free proof of Theorem~\ref{thm:main}:
\begin{itemize}
\item \texttt{period\_even}: for \emph{every} $p\ge 3$ (primality of $p$ is
not needed) every period $n$ is even and $\ne 2$;
\item \texttt{only\_two}: a prime period pair $(p,n)$ of primes forces
$p=2$;
\item \texttt{pisano2}, \dots, \texttt{pisano13}: the referee's numbers
$\pi(2){=}3$, $\pi(3){=}8$, $\pi(5){=}20$, $\pi(7){=}16$, $\pi(11){=}10$,
$\pi(13){=}28$, each a single \texttt{rfl} computation that also verifies
minimality and the wrap-around.
\end{itemize}
Every declaration is audited by \texttt{Check.lean} with
\texttt{\#print axioms} and depends on \emph{no} axioms.  To keep the audit
clean the development avoids \texttt{omega}, \texttt{simp} and
\texttt{decide} on symbolic goals; all congruences are phrased
existentially.

\section{Independent recomputation}

\texttt{reproduce.py} recomputes $\pi(p)$ by fast doubling for the
referee's data points, verifies that the only prime below $2000$ with
prime Pisano period is $2$, and checks that all other primes below $2000$
have even $\pi(p)$.

\section{Scope}

The disproof is unconditional: no variant of the conjecture survives.
\end{document}
